% Diapositiva 13

% Luis 
%============================================
% 11. Explicar Procedimientos para agregar ruido Gaussiano
%=============================================

\begin{frame}{Cálculo Paso a Paso: Distancia Euclidiana ($R_1$)}

    \begin{columns}[T]
        %========================================
        % COLUMNA IZQUIERDA: Contexto (Reducida al 30%)
        %========================================
        \column{0.30\textwidth}
        \begin{block}{Datos del Sensor 1}
            \scriptsize
            \textbf{S1:} $(-4.0, -3.0, 0.2)$ km \\
            \textbf{Fuente:} $(x_{\text{real}}, y_{\text{real}}, z_{\text{real}}) = (1.2, -0.8, -1.1)$ km
        \end{block}

        \vspace{0.4cm}

        \begin{block}{Nota Técnica}
            \tiny
            Sustitución de parámetros en la métrica euclidiana para determinar el radio de propagación.
        \end{block}

        %========================================
        % COLUMNA DERECHA: Operación (Ampliada al 68%)
        %========================================
        \column{0.68\textwidth}
        \begin{exampleblock}{Procedimiento Detallado}
            \centering
            \footnotesize % Tamaño ideal para que la raíz no se corte
            \begin{align*}
                R_1 &= \sqrt{(x_1 - x_{\text{real}})^2 + (y_1 - y_{\text{real}})^2 + (z_1 - z_{\text{real}})^2} \\[0.15cm]
                R_1 &= \sqrt{(-4.0 - 1.2)^2 + (-3.0 - (-0.8))^2 + (0.2 - (-1.1))^2} \\[0.15cm]
                R_1 &= \sqrt{(-5.2)^2 + (-3.0 + 0.8)^2 + (0.2 + 1.1)^2} \\[0.15cm]
                R_1 &= \sqrt{(-5.2)^2 + (-2.2)^2 + (1.3)^2} \\[0.15cm]
                R_1 &= \sqrt{27.04 + 4.84 + 1.69} \\[0.15cm]
                R_1 &= \sqrt{33.57} \implies \mathbf{R_1 \approx 5.794}
            \end{align*}
        \end{exampleblock}
    \end{columns}

    \vspace{0.3cm}
    \begin{center}
        \tiny \textcolor{gray}{Cálculo fundamental para la posterior aplicación del modelo de atenuación.}
    \end{center}

\end{frame}

